% =========================================================
% Oscillation --- Notes
% Chapter: continuity
% Volume III - Analysis
% =========================================================

\subsection{Oscillation}
\label{sec:oscillation}

\begin{tcolorbox}[title={Toolkit --- Oscillation}]
\begin{itemize}
  \item $\operatorname{OscillationOn}(f,E,\omega)$:
    $\omega(f;E)=\sup_{x_1,x_2\in E}|f(x_1)-f(x_2)|=\operatorname{diam}(f(E))$.
  \item $\operatorname{OscillationAt}(f,x_0,\omega)$:
    $\omega(f;x_0)=\lim_{\delta\to 0^+}\omega(f;U_\delta(x_0)\cap E)$.
    The limit exists because the function $\delta\mapsto\omega(f;U_\delta(x_0)\cap E)$
    is nonincreasing.
  \item $\operatorname{ContinuousAt}(f,x_0)\iff\omega(f;x_0)=0$.
  \item Discontinuity set: $D(f)=\{x:\omega(f;x)>0\}=
    \bigcup_{n=1}^\infty\{x:\omega(f;x)\geq 1/n\}$.
  \item Each threshold set is closed; $D(f)$ is $F_\sigma$.
\end{itemize}
\end{tcolorbox}

% ---------------------------------------------------------
% Definition --- Oscillation on a Set
% ---------------------------------------------------------

\begin{tcolorbox}[colback=defbox, colframe=defborder, arc=2pt,
  left=6pt, right=6pt, top=4pt, bottom=4pt,
  title={\small\textbf{Definition (Oscillation on a Set)}},
  fonttitle=\small\bfseries]
\begin{definition}[Oscillation on a Set]
\label{def:oscillation-on-a-set}
Let $A\subseteq\mathbb{R}$, let $f:A\to\mathbb{R}$, and let
$E\subseteq A$. The oscillation of $f$ on $E$ is the extended real number
\[
\omega(f;E)\coloneqq
\sup\{\,|f(x_1)-f(x_2)|:x_1,x_2\in E\,\}
=\mathrm{diam}(f(E)).
\]
\end{definition}
\end{tcolorbox}

\begin{remark*}[Standard quantified statement]
\[
\begin{aligned}
&\text{Let } A\subseteq\mathbb{R},\; f:A\to\mathbb{R},\; E\subseteq A,\;
\omega\in[0,\infty].\\
&\omega(f;E)=\omega
\Longleftrightarrow
\omega=\sup\{\,|f(x_1)-f(x_2)|:x_1,x_2\in E\,\}.
\end{aligned}
\]
\end{remark*}

\begin{remark*}[Definition predicate reading]
\[
\operatorname{OscillationOn}(f,E,\omega)
\coloneqq
\omega=\sup\{\,|f(x_1)-f(x_2)|:x_1,x_2\in E\,\}.
\]
\end{remark*}

\begin{remark*}[Negated quantified statement]
\[
\begin{aligned}
&\text{Let } A\subseteq\mathbb{R},\; f:A\to\mathbb{R},\; E\subseteq A,\;
\omega\in[0,\infty].\\
&\omega(f;E)\ne\omega
\Longleftrightarrow
\omega\ne\sup\{\,|f(x_1)-f(x_2)|:x_1,x_2\in E\,\}.
\end{aligned}
\]
\end{remark*}

\begin{remark*}[Negation predicate reading]
\[
\neg\operatorname{OscillationOn}(f,E,\omega)
\Longleftrightarrow
\omega\ne\sup\{\,|f(x_1)-f(x_2)|:x_1,x_2\in E\,\}.
\]
\end{remark*}

\begin{remark*}[Failure modes]
The proposed value $\omega$ fails either because it is not an upper
bound for pairwise differences of function values on $E$, or because it
is an upper bound but not the least such upper bound.
\end{remark*}

\begin{remark*}[Failure mode decomposition]
\[
\underbrace{\exists x_1,x_2\in E:\; |f(x_1)-f(x_2)|>\omega}_{\text{not an upper bound}}
\;\lor\;
\underbrace{\exists \eta<\omega\; \forall x_1,x_2\in E:\; |f(x_1)-f(x_2)|\le \eta}_{\text{not least}}.
\]
\end{remark*}

\begin{remark*}[Interpretation]
Oscillation on a set records the spread of the image --- the diameter
of $f(E)$. Small oscillation means all values of $f$ on $E$ lie
close together. Zero oscillation means $f$ is constant on $E$.
\end{remark*}

\NoLocalDependencies
% ---------------------------------------------------------
% Definition --- Oscillation at a Point
% ---------------------------------------------------------

\begin{tcolorbox}[colback=defbox, colframe=defborder, arc=2pt,
  left=6pt, right=6pt, top=4pt, bottom=4pt,
  title={\small\textbf{Definition (Oscillation at a Point)}},
  fonttitle=\small\bfseries]
\begin{definition}[Oscillation at a Point]
\label{def:oscillation-at-a-point}
Let $E \subseteq \mathbb{R}$, let $f:E \to \mathbb{R}$, and let $x_0 \in \mathbb{R}$ satisfy $U_\delta(x_0)\cap E \neq \varnothing$ for every $\delta>0$, where $U_\delta(x_0)=\{x \in \mathbb{R}: |x-x_0|<\delta\}$. The oscillation of $f$ at $x_0$ is the extended real number
\[
\omega(f;x_0)
=
\lim_{\delta \to 0^+}
\omega\bigl(f;\,U_\delta(x_0)\cap E\bigr),
\]
with $\omega(f;A)$ denoting the oscillation of $f$ on the set $A$. Because the function $\delta \mapsto \omega\bigl(f;\,U_\delta(x_0)\cap E\bigr)$ is nonincreasing, this right-hand limit exists (possibly $+\infty$).
\end{definition}
\end{tcolorbox}

\begin{remark*}[Standard quantified statement]
\[
\begin{aligned}
&\text{Let } E \subseteq \mathbb{R},\, f:E \to \mathbb{R},\, x_0 \in \mathbb{R},\, \omega \in [0,\infty].\\
&\bigl[(\forall \delta>0:\, U_\delta(x_0)\cap E \neq \varnothing)\ \land\ \omega(f;x_0)=\omega\bigr]\\
&\Longleftrightarrow
\forall \varepsilon>0\; \exists \delta>0\; \forall r\;(0<r<\delta \Rightarrow |\omega(f;U_r(x_0)\cap E)-\omega|<\varepsilon).
\end{aligned}
\]
\end{remark*}

\begin{remark*}[Definition predicate reading]
\[
\operatorname{OscillationAt}(f,x_0,\omega) \coloneqq
\lim_{\delta \to 0^+} \omega\bigl(f;\,U_\delta(x_0)\cap E\bigr) = \omega,
\quad U_\delta(x_0)=\{x \in \mathbb{R}: |x-x_0|<\delta\}.
\]
\end{remark*}

\begin{remark*}[Negated quantified statement]
\[
\begin{aligned}
&\text{Let } E \subseteq \mathbb{R},\, f:E \to \mathbb{R},\, x_0 \in \mathbb{R},\, \omega \in [0,\infty].\\
&\neg\bigl[(\forall \delta>0:\, U_\delta(x_0)\cap E \neq \varnothing)\ \land\ \omega(f;x_0)=\omega\bigr]\\
&\Longleftrightarrow
\Bigl[\exists \delta_0>0:\, U_{\delta_0}(x_0)\cap E = \varnothing\Bigr]\\
&\qquad\lor
\Bigl(\exists \varepsilon_0>0\; \forall \delta>0\; \exists r\;(0<r<\delta \land |\omega(f;U_r(x_0)\cap E)-\omega| \ge \varepsilon_0)\Bigr).
\end{aligned}
\]
\end{remark*}

\begin{remark*}[Negation predicate reading]
\[
\neg\operatorname{OscillationAt}(f,x_0,\omega)
\Longleftrightarrow
\Bigl[\exists \delta_0>0:\, U_{\delta_0}(x_0)\cap E = \varnothing\Bigr]
\lor
\Bigl(\exists \varepsilon_0>0\; \forall \delta>0\; \exists r\;(0<r<\delta \land |\omega(f;U_r(x_0)\cap E)-\omega| \ge \varepsilon_0)\Bigr).
\]
\end{remark*}

\begin{remark*}[Failure modes]
The definition fails in two distinct ways: either the point $x_0$ has a punctured neighborhood missing $E$, so there are no arbitrarily small samples to measure, or the local oscillation values near $x_0$ refuse to stabilize, preventing the limit from existing.
\end{remark*}

\begin{remark*}[Failure mode decomposition]
\[
\neg\operatorname{OscillationAt}(f,x_0,\omega)
=
\underbrace{\exists \delta_0>0:\, U_{\delta_0}(x_0)\cap E = \varnothing}_{\text{no nearby sample points}}
\;\lor\;
\underbrace{\exists \varepsilon_0>0\; \forall \delta>0\; \exists r\;(0<r<\delta \land |\omega(f;U_r(x_0)\cap E)-\omega| \ge \varepsilon_0)}_{\text{oscillation limit does not exist}}.
\]
\end{remark*}

\begin{remark*}[Interpretation]
Oscillation at $x_0$ measures the largest possible jump in function values that survives after one zooms in arbitrarily tightly around $x_0$. The shrinking neighborhoods $U_r(x_0)\cap E$ supply finer and finer samples, and the oscillation on each such set records the spread of $f$ there. If these spreads settle to a single number, that number captures the local variability; it is zero precisely when $f$ becomes arbitrarily close to a single value, i.e., when $f$ is continuous at $x_0$. Failure occurs either because $x_0$ has no nearby domain points or because the spreads fluctuate indefinitely, reflecting wild local behavior.
\end{remark*}

\begin{remark*}[Dependencies]
\hyperref[def:oscillation-on-a-set]{Definition~\ref*{def:oscillation-on-a-set}}
\end{remark*}


% ---------------------------------------------------------
% Theorem --- Continuity iff Zero Oscillation
% ---------------------------------------------------------

\begin{tcolorbox}[colback=thmbox, colframe=thmborder, arc=2pt,
  left=6pt, right=6pt, top=4pt, bottom=4pt,
  title={\small\textbf{Theorem (Continuity Characterised by Oscillation)}},
  fonttitle=\small\bfseries]
\begin{theorem}[Continuity Characterised by Oscillation]
\label{thm:continuity-iff-zero-oscillation}
Let $A \subseteq \mathbb{R}$, let $f \colon A \to \mathbb{R}$, and let $c \in A$. Then $f$ is continuous at $c$ if and only if $\omega(f;c)=0$.
\hyperref[prf:continuity-iff-zero-oscillation]{Go to proof.}
\end{theorem}
\end{tcolorbox}

\begin{remark*}[Standard quantified statement]
\[
\begin{aligned}
&\text{Let } A \subseteq \mathbb{R},\; f \colon A \to \mathbb{R},\; c \in A.\\
&\left(\forall \varepsilon>0 \; \exists \delta>0 \; \forall x \in A \; \bigl(|x-c|<\delta \Rightarrow |f(x)-f(c)|<\varepsilon\bigr)\right)\\
&\Longleftrightarrow\\
&\left(\inf_{\delta>0} \sup \bigl\{|f(x)-f(y)| : x,y \in A,\; |x-c|<\delta,\; |y-c|<\delta \bigr\} = 0\right).
\end{aligned}
\]
\end{remark*}

\begin{remark*}[Predicate reading]
\[
\operatorname{ContinuousAt}(f,c) \Longleftrightarrow \operatorname{OscillationAt}(f,c,0).
\]
\end{remark*}

\begin{remark*}[Negated quantified statement]
\[
\begin{aligned}
&\text{Let } A \subseteq \mathbb{R},\; f \colon A \to \mathbb{R},\; c \in A.\\
&\Bigl(\forall \varepsilon>0 \; \exists \delta>0 \; \forall x \in A \; (|x-c|<\delta \Rightarrow |f(x)-f(c)|<\varepsilon)\Bigr)
\land \Bigl(\inf_{\delta>0}\sup\{|f(x)-f(y)| : x,y \in A,\; |x-c|<\delta,\; |y-c|<\delta\}>0\Bigr)\\
&\qquad \lor\\
&\Bigl(\exists \varepsilon_0>0 \; \forall \delta>0 \; \exists x \in A \; (|x-c|<\delta \land |f(x)-f(c)| \ge \varepsilon_0)\Bigr)
\land \Bigl(\inf_{\delta>0}\sup\{|f(x)-f(y)| : x,y \in A,\; |x-c|<\delta,\; |y-c|<\delta\}=0\Bigr).
\end{aligned}
\]
\end{remark*}

\begin{remark*}[Negation predicate reading]
\[
\neg\bigl(\operatorname{ContinuousAt}(f,c) \Longleftrightarrow \operatorname{OscillationAt}(f,c,0)\bigr).
\]
\end{remark*}

\begin{remark*}[Failure modes]
The equivalence fails precisely in two mutually exclusive ways: either $f$ remains continuous at $c$ while the oscillation retains a positive lower bound, or $\omega(f;c)$ vanishes even though $f$ does not satisfy the $\varepsilon$-$\delta$ continuity condition at $c$.
\end{remark*}

\begin{remark*}[Failure mode decomposition]
\[
\underbrace{\operatorname{ContinuousAt}(f,c) \land \neg \operatorname{OscillationAt}(f,c,0)}_{\text{Positive oscillation despite continuity}}
\;\lor\;
\underbrace{\neg \operatorname{ContinuousAt}(f,c) \land \operatorname{OscillationAt}(f,c,0)}_{\text{Zero oscillation without continuity}}.
\]
\end{remark*}

\begin{remark*}[Interpretation]
The theorem states that measuring the oscillation of $f$ near $c$ captures exactly the same local regularity as the classical $\varepsilon$-$\delta$ definition of continuity. If oscillation collapses to zero, every pair of nearby function values becomes arbitrarily close, forcing $f(x)$ to approach $f(c)$. Conversely, any discontinuity manifests as a nonzero oscillation, because nearby points can produce separated outputs. The standard failure therefore arises from detecting a residual jump or oscillatory behavior in arbitrarily small neighborhoods of $c$.
\end{remark*}

\begin{remark*}[Dependencies]
\hyperref[def:continuous-at-point]{Definition~\ref*{def:continuous-at-point}},
\hyperref[def:oscillation-at-a-point]{Definition~\ref*{def:oscillation-at-a-point}}.
\end{remark*}


% ---------------------------------------------------------
% Proposition --- Discontinuity Set via Oscillation
% ---------------------------------------------------------

\begin{proposition}[Discontinuity Set via Oscillation]
\label{prop:discontinuity-set-via-oscillation}
Let $f:E\to\mathbb{R}$. Then
\[
D(f):=\{x\in E: f \text{ is discontinuous at } x\}
=\{x\in E: \omega(f;x)>0\}
=\bigcup_{n=1}^{\infty} \{x\in E: \omega(f;x)\ge 1/n\}.
\]
\hyperref[prf:discontinuity-set-via-oscillation]{Go to proof.}
\end{proposition}

\begin{remark*}[Standard quantified statement]
\[
\begin{aligned}
&\text{Let } f:E\to\mathbb{R}.\\
&\forall x\in E:\;
x\in D(f)
\Longleftrightarrow \omega(f;x)>0
\Longleftrightarrow
\exists n\in\mathbb{N}:\; \omega(f;x)\ge \tfrac1n.
\end{aligned}
\]
\end{remark*}

\begin{remark*}[Predicate reading]
\[
\begin{aligned}
&x\in D(f)
\Longleftrightarrow
\operatorname{DiscontinuousAt}(f,x)\\
&\Longleftrightarrow
\exists \omega\;(\operatorname{OscillationAt}(f,x,\omega)\land \omega>0)\\
&\Longleftrightarrow
\exists n\in\mathbb{N}\;\exists \omega\;(\operatorname{OscillationAt}(f,x,\omega)\land \omega\ge \tfrac1n).
\end{aligned}
\]
\end{remark*}

\begin{remark*}[Negated quantified statement]
\[
\begin{aligned}
&\text{Let } f:E\to\mathbb{R}.\\
&\neg\Bigl[\forall x\in E\;\bigl(x\in D(f)\Longleftrightarrow \omega(f;x)>0 \Longleftrightarrow \exists n\in\mathbb{N}\; \omega(f;x)\ge \tfrac1n\bigr)\Bigr]\\
&\Longleftrightarrow \exists x\in E\;\Bigl((x\in D(f)\land \omega(f;x)=0)\lor(x\notin D(f)\land \omega(f;x)>0)\\
&\qquad\qquad\lor(\omega(f;x)>0\land \forall n\in\mathbb{N}\; \omega(f;x)<\tfrac1n)\Bigr).
\end{aligned}
\]
\end{remark*}

\begin{remark*}[Negation predicate reading]
\[
\begin{aligned}
\exists x\in E\;\Bigl(
&\exists \omega\;(\operatorname{DiscontinuousAt}(f,x)\land \operatorname{OscillationAt}(f,x,\omega)\land \omega=0)\;\lor\\
&\exists \omega\;(\neg\operatorname{DiscontinuousAt}(f,x)\land \operatorname{OscillationAt}(f,x,\omega)\land \omega>0)\;\lor\\
&\exists \omega\;(\operatorname{OscillationAt}(f,x,\omega)\land \omega>0\land \forall n\in\mathbb{N}\; \omega<\tfrac1n)
\Bigr).
\end{aligned}
\]
\end{remark*}

\begin{remark*}[Failure modes]
The proposition fails precisely when some $x\in E$ is marked as discontinuous although its oscillation vanishes, or when $x$ is declared continuous despite a positive oscillation, or when a positive oscillation never exceeds any reciprocal integer, contradicting the Archimedean step that converts $\omega(f;x)>0$ into the countable union description.
\end{remark*}

\begin{remark*}[Failure mode decomposition]
\[
\exists x\in E\;
\begin{cases}
x\in D(f)\land \omega(f;x)=0 & \text{(Type I)},\\[4pt]
x\notin D(f)\land \omega(f;x)>0 & \text{(Type II)},\\[4pt]
\omega(f;x)>0\land \forall n\in\mathbb{N}\; \omega(f;x)<\tfrac1n & \text{(Type III)}.
\end{cases}
\]
\end{remark*}

\begin{remark*}[Interpretation]
Oscillation detects discontinuity quantitatively: every discontinuity forces $\omega(f;x)$ to stay above some positive threshold, and conversely the vanishing of oscillation characterises continuity. Because each set $\{x\in E:\omega(f;x)\ge 1/n\}$ is closed, the discontinuity set $D(f)$ is an $F_{\sigma}$, a countable union of closed sets. This structural description is the gateway to measure-theoretic criteria such as the Lebesgue characterisation of Riemann integrability, where the size of $D(f)$ controls integrability.
\end{remark*}

\begin{remark*}[Dependencies]
\hyperref[def:point-of-discontinuity]{Definition~\ref*{def:point-of-discontinuity}};
\hyperref[def:oscillation-at-a-point]{Definition~\ref*{def:oscillation-at-a-point}}.
\end{remark*}
